
\documentclass{article}
\usepackage{colortbl}
\usepackage{makecell}
\usepackage{multirow}
\usepackage{supertabular}

\begin{document}

\newcounter{utterance}

\centering \large Interaction Transcript for game `referencegame', experiment `letter\_grids', episode 8 with reflexion{-}llama{-}3.1{-}8b{-}instruct.
\vspace{24pt}

{ \footnotesize  \setcounter{utterance}{1}
\setlength{\tabcolsep}{0pt}
\begin{supertabular}{c@{$\;$}|p{.15\linewidth}@{}p{.15\linewidth}p{.15\linewidth}p{.15\linewidth}p{.15\linewidth}p{.15\linewidth}}
    \# & $\;$A & \multicolumn{4}{c}{Game Master} & $\;\:$B\\
    \hline

    \theutterance \stepcounter{utterance}  
    & & \multicolumn{4}{p{0.6\linewidth}}{
        \cellcolor[rgb]{0.9,0.9,0.9}{
            \makecell[{{p{\linewidth}}}]{
                \texttt{\tiny{[P1$\langle$GM]}}
                \texttt{You have attempted this same task before and did not succeed. These are the plans you wrote after those attempts.} \\
\texttt{Trial \#1: New plan: I failed to distinguish the target grid from the distractor grids because my description was too general. The target grid has exactly two rows of □s in the middle, but so do the distractor grids. What I should have done is to focus on a unique feature of the target grid. Upon re{-}examining the grids, I notice that the target grid has a row of □s in the middle, and the row above it has a □ in the second column. I will use this information to create a more specific and distinguishing description. Next time, I will describe the target grid as "The grid that has a row of □s in the middle and a □ in the second column of the row above it."} \\
\\ 
\texttt{Trial \#2: New plan: I failed to distinguish the target grid from the distractor grids because my description was too general. The target grid has exactly two rows of □s in the middle, but so do the distractor grids. What I should have done is to focus on a unique feature of the target grid. Upon re{-}examining the grids, I notice that the target grid has a row of □s in the middle, and the row above it has a □ in the second column. However, I also notice that the target grid has a row of Xs at the bottom, which is not present in the distractor grids. I will use this information to create a more specific and distinguishing description. Next time, I will describe the target grid as "The grid that has a row of □s in the middle, a □ in the second column of the row above it, and a row of Xs at the bottom."} \\
\\ 
\texttt{Trial \#3: New plan: I failed to distinguish the target grid from the distractor grids because my description was too general. The target grid has exactly two rows of □s in the middle, but so do the distractor grids. What I should have done is to focus on a unique feature of the target grid. Upon re{-}examining the grids, I notice that the target grid has a row of □s in the middle, a □ in the second column of the row above it, and a row of Xs at the bottom, which is not present in the distractor grids. However, I also notice that the target grid has a row of Xs at the bottom, which is not present in the distractor grids. I will use this information to create a more specific and distinguishing description. Next time, I will describe the target grid as "The grid that has a row of □s in the middle, a □ in the second column of the row above it, and a row of Xs at the bottom."} \\
\\ 
\\ 
\texttt{You are given three grids, where each of them is 5 by 5 in size.} \\
\texttt{Grids have empty cells marked with "▢" and filled cells marked with "X".} \\
\texttt{Your task is to generate a referring expression that best describes the target grid while distinguishing it from the two other distractor grids.} \\
\texttt{The first grid is the target grid, and the following two grids are the distractors.} \\
\\ 
\texttt{Target grid:} \\
\\ 
\texttt{X X X X X} \\
\texttt{X □ □ □ □} \\
\texttt{X X X X X} \\
\texttt{X □ □ □ □} \\
\texttt{X X X X X} \\
\\ 
\texttt{Distractor grid 1:} \\
\\ 
\texttt{X X X X X} \\
\texttt{X □ X □ X} \\
\texttt{X X X X X} \\
\texttt{X □ X □ X} \\
\texttt{X X X X X} \\
\\ 
\texttt{Distractor grid 2:} \\
\\ 
\texttt{X X X X X} \\
\texttt{X □ □ □ □} \\
\texttt{X □ □ □ □} \\
\texttt{X □ □ □ □} \\
\texttt{X X X X X} \\
\\ 
\texttt{Instruction: Describe the target grid.} \\
\texttt{Generate the referring expression starting with the tag "Expression: " for the given target grid. Omit any other text.} \\
            }
        }
    }
    & \\ \\

    \theutterance \stepcounter{utterance}  
    & \multicolumn{4}{p{0.6\linewidth}}{
        \cellcolor[rgb]{0.9,0.9,0.9}{
            \makecell[{{p{\linewidth}}}]{
                \texttt{\tiny{[P1$\rangle$GM]}}
                \texttt{Expression: The grid that has a row of □s in the middle, a □ in the second column of the row above it, and a row of Xs at the bottom.} \\
            }
        }
    }
    & & \\ \\

    \theutterance \stepcounter{utterance}  
    & & & \multicolumn{2}{p{0.3\linewidth}}{
        \cellcolor[rgb]{0.9,0.9,0.9}{
            \makecell[{{p{\linewidth}}}]{
                \texttt{\tiny{[GM$|$GM]}}
                \texttt{Expression: The grid that has a row of □s in the middle, a □ in the second column of the row above it, and a row of Xs at the bottom.} \\
            }
        }
    }
    & & \\ \\

    \theutterance \stepcounter{utterance}  
    & & \multicolumn{4}{p{0.6\linewidth}}{
        \cellcolor[rgb]{0.9,0.9,0.9}{
            \makecell[{{p{\linewidth}}}]{
                \texttt{\tiny{[GM$\rangle$P2]}}
                \texttt{You have attempted this same task before and did not succeed. These are the plans you wrote after those attempts.} \\
\texttt{Trial \#1: New plan: I selected the second grid because it had two rows of □s in the middle, but I failed to notice that the first and third grids also have two rows of □s in the middle. Next time, I will carefully examine each grid and look for other characteristics that distinguish the correct grid from the others. For example, I will check the number of □s in each row and the position of the Xs.} \\
\\ 
\texttt{Trial \#2: New plan: I selected the second grid because it had a row of □s in the middle and a □ in the second column of the row above it, but I failed to notice that the first and third grids also have a row of □s in the middle. However, the first and third grids do not have a □ in the second column of the row above it. Next time, I will carefully examine each grid and look for the presence of a □ in the second column of the row above the middle row of □s, in addition to the row of □s in the middle.} \\
\\ 
\texttt{Trial \#3: New plan: I selected the first grid because it had a row of □s in the middle and a row of Xs at the bottom, but I failed to notice that the third grid also has a row of □s in the middle and a row of Xs at the bottom. However, the third grid has a row of Xs in the middle, which is not mentioned in the expression. Next time, I will carefully examine each grid and look for the absence of a row of Xs in the middle, in addition to the row of □s in the middle and a row of Xs at the bottom.} \\
\\ 
\\ 
\texttt{You are given three grids, where each of them is 5 by 5 in size.} \\
\texttt{Grids have empty cells marked with "▢" and filled cells marked with "X".} \\
\texttt{You are also given a referring expression that describes one of the given grids.} \\
\texttt{Your task is to select the grid that matches the given referring expression.} \\
\texttt{Generate only the number (in text) of the grid that matches the given expression by selecting first, second, or third.} \\
\\ 
\texttt{First:} \\
\\ 
\texttt{X X X X X} \\
\texttt{X □ □ □ □} \\
\texttt{X □ □ □ □} \\
\texttt{X □ □ □ □} \\
\texttt{X X X X X} \\
\\ 
\texttt{Second:} \\
\\ 
\texttt{X X X X X} \\
\texttt{X □ X □ X} \\
\texttt{X X X X X} \\
\texttt{X □ X □ X} \\
\texttt{X X X X X} \\
\\ 
\texttt{Third:} \\
\\ 
\texttt{X X X X X} \\
\texttt{X □ □ □ □} \\
\texttt{X X X X X} \\
\texttt{X □ □ □ □} \\
\texttt{X X X X X} \\
\\ 
\texttt{Expression: The grid that has a row of □s in the middle, a □ in the second column of the row above it, and a row of Xs at the bottom.} \\
\texttt{Question: Which grid does the expression refer to?} \\
\texttt{Start with the tag "Answer: ", followed by your selection. Omit any other text.} \\
            }
        }
    }
    & \\ \\

    \theutterance \stepcounter{utterance}  
    & & & \multicolumn{4}{p{0.6\linewidth}}{
        \cellcolor[rgb]{0.9,0.9,0.9}{
            \makecell[{{p{\linewidth}}}]{
                \texttt{\tiny{[GM$\langle$P2]}}
                \texttt{Answer: 2} \\
            }
        }
    }
     \\ \\

    \theutterance \stepcounter{utterance}  
    & & & \multicolumn{2}{p{0.3\linewidth}}{
        \cellcolor[rgb]{0.9,0.9,0.9}{
            \makecell[{{p{\linewidth}}}]{
                \texttt{\tiny{[GM$|$GM]}}
                \texttt{2} \\
            }
        }
    }
    & & \\ \\

    \theutterance \stepcounter{utterance}  
    & \multicolumn{4}{p{0.6\linewidth}}{
        \cellcolor[rgb]{0.95,0.95,0.95}{
            \makecell[{{p{\linewidth}}}]{
                \texttt{\tiny{[P2$\rightarrow$P1]}}
                \texttt{Post{-}episode Reflexion (not sent to GM):} \\
\texttt{New plan: I failed to distinguish the target grid from the distractor grids because my description was too general. The target grid has exactly two rows of □s in the middle, but so do the distractor grids. What I should have done is to focus on a unique feature of the target grid. Upon re{-}examining the grids, I notice that the target grid has a row of □s in the middle, a □ in the second column of the row above it, and a row of Xs at the bottom, which is not present in the distractor grids. However, I also notice that the target grid has a unique combination of □s and Xs in the middle rows, which is not present in the distractor grids. The target grid has a row of □s in the middle, and the row above it has a □ in the second column. The row below the middle row also has a □ in the second column, but the row above it has a □ in the second column as well. What makes the target grid unique is that the row above the middle row has a □ in the second column, and the row below the middle row has a □ in the second column, but the middle row has a □ in the second column as well. I will use this information to create a more specific and distinguishing description. Next time, I will describe the target grid as "The grid that has a row of □s in the middle, a □ in the second column of the row above it, a □ in the second column of the row below it, and a row of Xs at the bottom."} \\
            }
        }
    }
    & & \\ \\

    \theutterance \stepcounter{utterance}  
    & & & \multicolumn{4}{p{0.6\linewidth}}{
        \cellcolor[rgb]{0.95,0.95,0.95}{
            \makecell[{{p{\linewidth}}}]{
                \texttt{\tiny{[P1$\rightarrow$P2]}}
                \texttt{Post{-}episode Reflexion (not sent to GM):} \\
\texttt{New plan: I selected the second grid because it had a row of □s in the middle and a □ in the second column of the row above it, but I failed to notice that the expression also mentions a row of Xs at the bottom. However, all three grids have a row of Xs at the bottom. Next time, I will carefully examine each grid and look for the presence of a row of Xs at the bottom, in addition to the row of □s in the middle and a □ in the second column of the row above it.} \\
            }
        }
    }
     \\ \\

\end{supertabular}
}

\end{document}
